\chapter{圆锥曲线}

%% ---- 知识框架 ---- %%
\begin{knowledgeframe}
\begin{itemize}
	\item \textbf{直线与圆的位置关系}：$d<r$ 相交、$d=r$ 相切、$d>r$ 相离；弦长公式 $|AB|=2\sqrt{r^2-d^2}$。
	\item \textbf{抛物线}：$y^2=2px$ 的焦点 $F(\frac{p}{2},0)$，准线 $x=-\frac{p}{2}$；切线问题常转化为联立方程。
	\item \textbf{椭圆}：$\frac{x^2}{a^2}+\frac{y^2}{b^2}=1$（$a>b>0$），焦点在长轴上。
	\item \textbf{直线与圆锥曲线}：联立方程后用韦达定理处理弦长、中点、斜率等问题。
	\item \textbf{点与圆的位置关系}：圆上点的切线方程、割线与切线的转化。
\end{itemize}
\end{knowledgeframe}

%% ---- 第一节：直线与圆 ---- %%
\section{直线与圆}

\startexercise

\subsection*{A组}

\begin{choice}[source={浙江宁波镇海中学模拟}, topic={直线与位置关系}]
已知直线 $l: ax+by+c=0$，$A(x_1, y_1)$，$B(x_2, y_2)$（其中 $x_1 \ne x_2$），当 $\dfrac{ax_1+by_1+c}{ax_2+by_2+c} > 1$ 时，直线 $l$ 与直线 $AB$ 的位置关系为
\begin{tasks}(2)
	\task 垂直
	\task 平行
	\task 相交
	\task 以上位置关系都有可能
\end{tasks}
\end{choice}
\begin{choice-sol}
待补充解析。
\end{choice-sol}


\subsection*{B组}

\begin{choice}[source={2024年高考全国甲卷理科选择题\#12}, topic={等差中项与圆的弦长}]
已知 $b$ 是 $a,c$ 的等差中项，直线 $ax+by+c=0$ 与圆 $x^{2}+y^{2}+4y-1=0$ 交于 $A,B$ 两点，则 $\lvert AB\rvert$ 的最小值为
\begin{tasks}(4)
	\task $2$
	\task $3$
	\task $4$
	\task $2\sqrt{5}$
\end{tasks}
\end{choice}
\begin{choice-sol}
\textbf{答案 C}

待补充解析。
\end{choice-sol}


\subsection*{C组}

\begin{choice}[source={2021八省联考}, topic={抛物线的切线与割线}]
已知抛物线 $y^2=2px$ 上三点 $A(2,2),B,C$。直线 $AB,AC$ 是圆 $M: (x-2)^2+y^2=1$ 的两条切线，则直线 $BC$ 的方程为
\begin{tasks}(2)
	\task $x+2y+1=0$
	\task $3x+6y+4=0$
	\task $2x+6y+3=0$
	\task $x+3y+2=0$
\end{tasks}
\end{choice}
\begin{choice-sol}
\textbf{答案 A}

待补充解析。
\end{choice-sol}


\stopexercise

%% ---- 第二节：椭圆 ---- %%
\section{椭圆}

\startexercise

\subsection*{习题集}

\begin{prob}[source={2011大纲卷理}, topic={椭圆}, score=12]
已知 $O$ 为坐标原点，$F$ 为椭圆 $C: x^2 + \dfrac{y^2}{2} = 1$ 在 $y$ 轴正半轴上的焦点，过 $F$ 且斜率为 $-\sqrt{2}$ 的直线 $l$ 与 $C$ 交于 $A, B$ 两点，点 $P$ 满足 $\overrightarrow{OA} + \overrightarrow{OB} + \overrightarrow{OP} = \mathbf{0}$。
\begin{enumerate}
	\item[(1)] 证明：点 $P$ 在 $C$ 上；
	\item[(2)] 设点 $P$ 关于点 $O$ 的对称点为 $Q$，证明：$A, P, B, Q$ 四点共圆。
\end{enumerate}
\end{prob}
\begin{prob-sol}
待补充解析。
\end{prob-sol}


\stopexercise
